\documentclass[10pt,a4paper]{amsart}
\usepackage[margin=19mm]{geometry}
\usepackage{amsmath,amssymb,mathtools}
\usepackage[scr=rsfs,cal=euler]{mathalfa}
\usepackage{xfrac}
\usepackage[unicode,hidelinks,bookmarks=false]{hyperref}
\newcommand{\ie}{i.e.\ }
\newcommand{\cf}{cf.\ }
\newcommand{\resp}{respectively\ }
\newcommand{\Subsection}{\subsection}
\providecommand{\nobreakdash}{\nobreak\hbox{-}}
\setlength{\emergencystretch}{2em}
\multlinegap0pt
\usepackage{bm}
\usepackage{bbm}
\usepackage{dsfont}
\usepackage{xspace}
\usepackage{cancel}
\usepackage{mathtools}
\usepackage{amsthm}
\newtheorem{theorem}{Theorem}
\newtheorem{lemma}[theorem]{Lemma}
\newtheorem{proposition}[theorem]{Proposition}
\newcommand{\R}{\mathbb R}
\newcommand{\Sph}{\mathbb S}
\newcommand{\Z}{\mathbb Z}
\title[Brownian Loops, Singular Homology, and Asymptotic Cycles]{Brownian Loops, Singular Homology, and Asymptotic Cycles}
\author{Alberto Verjovsky, Ricardo F. Vila-Freyer}
\date{Source: arXiv:2609.27851v1 (CC BY 4.0)}
\begin{document}
\maketitle

\noindent\emph{Source and licence.} Brownian Loops, Singular Homology, and Asymptotic Cycles. Version arXiv:2609.27851v1 (CC BY 4.0), distributed under the Creative Commons Attribution 4.0 International licence, as recorded in the arXiv licence metadata for this exact version and shown on the abstract page https://arxiv.org/abs/2609.27851v1. Full rights notice: licenses/arxiv-2609.27851-CC-BY-4.0.txt.\par\smallskip

\noindent\textit{Editorial scope.} Counted unit: the Proposition whose source label is \texttt{prop:pairing} in the article (source0.tex lines 329--338, source label \texttt{prop:pairing}), delivered together with the article's own complete original proof of it (source0.tex lines 340--358, one \texttt{proof} environment of 19 lines). No mathematics is written, repaired or re-derived: statement and proof are the article's own text line for line. Required context retained uncounted: lem:circlelift (lines 312--318). Every definition, display and label of those lines is retained unchanged. Omitted from this package and not redistributed: 1--311, 319--328, 339--339, 359--1112. Nothing inside a retained excerpt is omitted or altered: every retained body is delivered exactly as the article's lines are written, so no omission receipt of any kind accompanies this package. Citations: the retained excerpts carry no citation command at all, so no bibliography entry of the article is transcribed and no reference is redistributed. Reference adaptation: none was needed and none was made; every reference inside the delivered unit resolves inside it, so no unresolved reference or citation is produced. Printed slips kept literal: no printed word, symbol, hypothesis or number of the counted statement or of its proof was altered. Layout and class: the article's own class is \texttt{amsart} with packages including inputenc, babel, amsmath, amssymb, amsfonts, amsthm, microtype, enumitem, cite, url, xcolor, hyperref; the wrapper is the lane's pinned \texttt{amsart} editorial wrapper at 10pt on A4, which restates the theorem-like environments the retained text needs and adds the article's own macro definitions that the retained text needs (\texttt{\textbackslash R}, \texttt{\textbackslash Sph}, \texttt{\textbackslash Z}), with extra wrapper packages (bm, bbm, dsfont, xspace, cancel, mathtools). The delivered numbering is the unit's own. Assets: no retained span contains a figure environment, a graphics-inclusion command, a PDF-page inclusion, a table or a TikZ picture; no image file is copied into this package, the package declares no asset dependency and no image file is redistributed with it. Licence: version arXiv:2609.27851v1 of this article is distributed under the Creative Commons Attribution 4.0 International licence, as recorded in the arXiv licence metadata for this exact version and shown on the abstract page https://arxiv.org/abs/2609.27851v1. A full rights notice is delivered at licenses/arxiv-2609.27851-CC-BY-4.0.txt. Characters: every retained excerpt is pure ASCII, so the delivered engine, XeLaTeX with its default Latin Modern text font, reports no missing character.
\begin{lemma}[Circle-valued semimartingales and lifts]\label{lem:circlelift}
Let $\phi:M\to\Sph^1$ be smooth and let $X_s$, $0\le s\le t$, be a continuous semimartingale.  Put $Y_s=\phi(X_s)$ and let $\widetilde Y_s$ be any continuous lift of $Y_s$ to $\R$ under $u\mapsto e^{iu}$.  Then
\[
 \int_0^t \phi^*\!\left(\frac{d\theta}{2\pi}\right)\circ dX_s
 =\frac{\widetilde Y_t-\widetilde Y_0}{2\pi}.
\]
\end{lemma}

\begin{proposition}[Stochastic representation of the homology pairing]\label{prop:pairing}
Let $X$ be a continuous semimartingale on $M$.  For every smooth closed
$1$-form $\omega$,
\begin{equation}\label{eq:pairingidentity}
   I_t(\omega)=\langle[\omega],H_t\rangle
\end{equation}
almost surely, simultaneously for all $t\ge0$, after choosing the usual
continuous version of the Stratonovich integral.  In particular,
$I_t(\omega)$ depends only on the de Rham class of $\omega$.
\end{proposition}

\begin{proof}
First suppose that $\omega=df$.  By the Stratonovich chain rule,
\[
   \int_0^t df\circ dX_s=f(X_t)-f(X_0),
\]
while
\[
   \int_{\gamma_{X_t,X_0}}df=f(X_0)-f(X_t).
\]
Thus $I_t(df)=0$, so $I_t$ depends only on the de Rham class.

The image of $H^1(M;\Z)$ in $H^1(M;\R)$ is a full lattice and therefore spans $H^1(M;\R)$.  Let $\alpha\in H^1(M;\Z)$ and represent it by a smooth map $\phi:M\to\Sph^1$.  The form $\phi^*(d\theta/(2\pi))$ represents the real image of $\alpha$.  Lemma~\ref{lem:circlelift} identifies the stochastic integral over the Brownian part with the lifted angular increment of $\phi(X_s)$.  The ordinary integral over the closing path gives the remaining angular increment.  Their sum is therefore the total angular change of the closed loop $\phi\circ\Gamma(X,t)$, divided by $2\pi$, namely its degree.  Hence
\[
 I_t\!\left(\phi^*\frac{d\theta}{2\pi}\right)
 =\deg(\phi\circ\Gamma(X,t))
 =\langle\alpha,[\Gamma(X,t)]\rangle.
\]
The identity follows for integral classes and hence, by linearity, for all classes in $H^1(M;\R)$.
\end{proof}

\end{document}
